\documentclass[11pt]{article}
\usepackage[margin=1in]{geometry}
\usepackage{amsmath,amssymb,amsthm}
\usepackage{enumitem}
\usepackage{microtype}
\usepackage{hyperref}
\hypersetup{colorlinks=true,urlcolor=blue,linkcolor=black,citecolor=black}

\setlength{\parskip}{4pt plus 1pt}
\setlength{\parindent}{0pt}

\newtheorem{thm}{Theorem}
\newtheorem{lem}[thm]{Lemma}
\newtheorem{prop}[thm]{Proposition}
\newtheorem{cor}[thm]{Corollary}
\theoremstyle{definition}
\newtheorem{defn}[thm]{Definition}

\title{Layer-0 hidden-zero kills no triangulation\\
       \small{(short note for the locality program)}}
\date{}

\begin{document}

\maketitle

\begin{abstract}
For the planar $\mathrm{Tr}(\phi^3)$ ansatz at $n$ external legs, the
\emph{layer-$0$ kill} at the cyclic hidden-zero zone $\mathcal Z_{r,r+2}$
eliminates the coefficient $a_M$ of every chord multiset $M$ that, at
that zone, contains neither the special chord nor the bare and uses no
(companion, substitute) pair in full. We prove that no triangulation of
the $n$-gon is killed at any zone. The argument is two lines once the
right object is identified: the unique triangle of the triangulation
that contains the polygon edge $(r,r{+}1)$ has its third vertex $a$
falling into one of three exhaustive cases, each of which forces a
violation of the layer-$0$ kill criterion. No appeal to Meisters's
two-ears theorem is needed.
\end{abstract}

%--------------------------------------------------------------------------
\section{Setup}
%--------------------------------------------------------------------------

Let $n\ge 4$ and label the vertices $1,2,\dots,n$ cyclically. A
\emph{chord} is a pair $(i,j)$, $i\ne j$, that is not a polygon edge.
For $r\in\{1,\dots,n\}$, the zone $\mathcal Z_{r,r+2}$ is the
codimension-$2$ hidden-zero locus on which
\begin{equation}\label{eq:subs}
   X_{r+1,k} \;=\; X_{r,k} - X_{r,r+2}
   \qquad
   \text{for } k\notin\{r,r{+}1,r{+}2\}\text{ (cyclic),}
\end{equation}
together with the bare relation $X_{r+1,r-1} = -X_{r,r+2}$. Indices are
taken modulo $n$ throughout.

\begin{defn}[Step-1 / layer-$0$ kill]\label{defn:K}
A chord multiset $M$ of size $n{-}3$ is \emph{killable at
$\mathcal Z_{r,r+2}$}, written $M\in\mathcal K_{r,r+2}$, if both
\begin{itemize}[leftmargin=*,topsep=0pt,parsep=0pt,itemsep=2pt]
\item[\rm (K1)] $M$ contains neither the special $X_{r,r+2}$ nor the
       bare $X_{r+1,r-1}$;
\item[\rm (K2)] for every $k\in\{r{+}3,\dots,r{+}n{-}2\}$ (cyclic), $M$
       does not contain both members of the (companion, substitute) pair
       $(X_{r,k},\,X_{r+1,k})$.
\end{itemize}
If $M\in\mathcal K_{r,r+2}$ then a leading-Laurent argument with
$X_{r,r+2}\to\infty$ forces $a_M=0$. We do not use that argument here.
\end{defn}

\begin{defn}[Triangulation]
A \emph{triangulation} $T$ of the $n$-gon is a maximal set of pairwise
non-crossing chords; equivalently, a set of $n{-}3$ chords whose
addition to the polygon edges partitions the polygon into $n{-}2$
triangles.
\end{defn}

%--------------------------------------------------------------------------
\section{Result}
%--------------------------------------------------------------------------

\begin{thm}\label{thm:main}
For every $n\ge 4$, every triangulation $T$ of the $n$-gon, and every
$r\in\{1,\dots,n\}$,
\[
   T \;\notin\; \mathcal K_{r,r+2}.
\]
\end{thm}

\begin{proof}
Fix $r$ and let $T$ be any triangulation. The polygon edge $(r,r{+}1)$
is a side of exactly one triangle of $T$; call it $\Delta$, and let $a$
be its third vertex (so $a\notin\{r,r{+}1\}$). The two sides of $\Delta$
incident to $a$ are then $(r,a)$ and $(r{+}1,a)$.

We split on the value of $a$:

\smallskip\noindent
\textbf{Case 1: $a = r{+}2$.} Then the side $(r,a)=(r,r{+}2)$ is the
special chord, and it is a side of $\Delta$, hence is contained in $T$.
So the special $X_{r,r+2}\in T$, violating (K1). Therefore
$T\notin\mathcal K_{r,r+2}$.

\smallskip\noindent
\textbf{Case 2: $a = r{-}1$ (cyclic).} Then the side $(r{+}1,a) =
(r{+}1,r{-}1)$ is the bare chord, and it is a side of $\Delta$, hence
in $T$. So the bare $X_{r+1,r-1}\in T$, violating (K1). Therefore
$T\notin\mathcal K_{r,r+2}$.

\smallskip\noindent
\textbf{Case 3: $a\in\{r{+}3,\dots,r{+}n{-}2\}$ (cyclic).} Then the
sides $(r,a)$ and $(r{+}1,a)$ are both chords (neither is a polygon
edge, since $a$ is not adjacent to $r$ or $r{+}1$ on the polygon
boundary). They are sides of $\Delta$, hence both in $T$. But these are
exactly the (companion, substitute) pair
$(X_{r,a},\,X_{r+1,a})$, both contained in $T$. This violates (K2).
Therefore $T\notin\mathcal K_{r,r+2}$.

The three cases exhaust the possible values of $a$, so the conclusion
$T\notin\mathcal K_{r,r+2}$ holds. \qed
\end{proof}

\begin{cor}\label{cor:loc}
No layer-$0$ hidden-zero argument at any cyclic zone kills the
coefficient $a_T$ of any triangulation $T$.
\end{cor}

\begin{proof}
The layer-$0$ kill at $\mathcal Z_{r,r+2}$ acts only on coefficients
$a_M$ with $M\in\mathcal K_{r,r+2}$. By Theorem~\ref{thm:main}, $T$ is
in no such set. \qed
\end{proof}

\section*{Remarks}

\textbf{1.\ The geometric content.} The proof says: the polygon edge
$(r,r{+}1)$ sits inside exactly one triangle of $T$, and the third
vertex $a$ of that triangle determines which of (K1)'s two clauses or
(K2) is violated. This is purely local at the edge $(r,r{+}1)$ and uses
nothing about the rest of $T$.

\smallskip\noindent
\textbf{2.\ No two-ears theorem.} Earlier formulations of this lemma
appealed to Meisters's two-ears theorem (every triangulation of a
convex polygon with $n\ge 4$ vertices has at least two ears). The
two-ears theorem is unnecessary here: the polygon edge $(r,r{+}1)$ is a
boundary edge of \emph{any} triangulation, in particular always lies in
exactly one triangle, and that single fact does the work.

\smallskip\noindent
\textbf{3.\ What this does and does not give.} Theorem~\ref{thm:main}
guarantees triangulations survive Step-1 (layer-$0$). It says nothing
about Step-2 (the bare-relation equality chains), Step-3 (higher
Laurent layers), or whether non-triangulations are all eventually
killed. Those are separate questions; this lemma is the local-survival
half of the locality programme.

\end{document}
